\section{Introduction}

The financial crisis in 2007/08 and the following worldwide recession marked a turning point in the economic and public debate about optimal fiscal and monetary policy. In the aftermath of the great recession, economists discussed about the effectiveness of expansionary fiscal policy to stabilize the economy and stimulate the private sector. In this debate the fiscal multiplier has received special attention.

In the the aftermath of the crisis, governments and central banks launched unconventional and in part uncharted measures to mitigate the harmful consequences of a prolonged recession. The strong reduction in the nominal interest rate was very preeminent among those measures because many central banks reduced the interest rates down to zero. Since the zero lower bound was hit and monetary policy exhausted one of its main tools, fiscal policy became more relevant. This new state of affairs gave birth to a literature increasingly concerned with fiscal policy and the government spending multiplier. This literature stresses the effectiveness of government spending in cases where monetary policy has already brought interest rates to zero and the economy finds itself in a liquidity trap.
\par
\bigskip
In this context, the authors Christopher Erceg and Jesper Lindé address in the paper \textit{``Is there a Fiscal Free Lunch in a liquidity trap?''} the question whether policymakers should limit the size of the government spending. They do so by writing a model in which the liquidity trap is endogenously determined and then  by calculating the multiplier schedule for different calibrations. They also draw the very important distinction between marginal and average spending multiplier, which is crucial for policymakers to better understand the effects of changes in government spending. Finally, the paper provides new insights into how fiscal spending and austerity impact the government budget in a liquidity trap.
\par
\bigskip
In our paper, we try to replicate the basic results of \citet{Erceg.2014} using Dynare. It is organized as follows. In Section II, we discuss briefly the current literature about the government spending multiplier and explain the main results of the analyzed paper. In Section III, we concentrate on the key equations of the paper, explain the structure of the model and the underlying microeconomics of the simple New Keynesian model. In Section IV, we present the results of the impulse response functions to show how the model reacts to an exogenous consumption taste shock and a government spending hike. We also analyze how the results vary if we change the slope of the Phillips curve or the parameters of the monetary rule. In Section V, we dive deeper into the concept of multiplier and explain the distinction between \textit{marginal} and \textit{average} multiplier. We show that the variation of the fiscal multiplier depends highly on the size of the government response and on the average contract duration (the Calvo parameter). In Section VI, we analyze the impacts of an increase in government spending on the government budget and answer the question wether there is a fiscal free lunch in a liquidity trap. In Section VII, we provide some additional information on our workflow that could be useful for future participants of this seminar. In Section VIII, we conclude with final remarks.

